\subsection{Исследование влияния амплитудных и фазовых искажений}

В данном разделе оценивается влияние неидеальности формирования квадратурных составляющих цифровым гетеродином на динамический диапазон устройства. В соответствии с заданием, для анализа амплитудного и фазового дисбаланса использовался режим работы гетеродина \texttt{nco\_type = 'single'}. Квадратурные искажения приводят к нарушению строгой ортогональности I и Q каналов, следствием чего является неполное подавление зеркального канала при комплексной обработке и появление в спектре мощной побочной составляющей — зеркальной помехи.

На рисунке \ref{fig:sfdr_amv} представлена зависимость уровня SFDR от амплитудного дисбаланса (параметр \texttt{amv}) при нулевом фазовом искажении.

\begin{figure}[H]
    \centering
    \includegraphics[width=0.8\textwidth]{p4_a.PNG}
    \caption{Зависимость уровня подавления побочных составляющих (SFDR) от амплитудных искажений}
    \label{fig:sfdr_amv}
\end{figure}

График показывает нелинейное и весьма резкое падение SFDR при увеличении амплитудной асимметрии между каналами. Появление амплитудного дисбаланса всего в 0.02 приводит к деградации свободного от помех динамического диапазона с исходных $\approx 71$ дБн до $\approx 58$ дБн. При максимальном исследуемом значении \texttt{amv} = 0.2 значение SFDR опускается ниже 40 дБн.

На рисунке \ref{fig:sfdr_pmv} продемонстрирована зависимость SFDR от фазового дисбаланса (параметр \texttt{pmv}, измеряемый в градусах) при отсутствии амплитудных искажений.

\begin{figure}[H]
    \centering
    \includegraphics[width=0.8\textwidth]{p4_f.PNG}
    \caption{Зависимость уровня подавления побочных составляющих (SFDR) от фазовых искажений}
    \label{fig:sfdr_pmv}
\end{figure}

Фазовые ошибки оказывают еще более катастрофическое влияние на чистоту выходного спектра. Отклонение фазы квадратурного сигнала всего на $2^\circ$ от идеальных $90^\circ$ обрушивает значение SFDR с $\approx 71$ дБн до $\approx 35$ дБн. При фазовой ошибке в $20^\circ$ уровень зеркальной помехи возрастает настолько, что разница между амплитудой полезного сигнала и помехи сокращается до критических $\approx 15$ дБн.

\textbf{Выводы по результатам моделирования:}
Даже незначительные амплитудные и, в особенности, фазовые искажения (I/Q дисбаланс) в цифровом гетеродине приводят к стремительному росту амплитуды зеркальной составляющей. Это полностью нивелирует преимущества многоразрядных АЦП, критически ухудшая SFDR на выходе цифрового фильтра. Для обеспечения высокого динамического диапазона системы требуется строгий математический контроль ортогональности и равенства амплитуд опорных колебаний.
